\documentclass[11pt]{article}

\usepackage[margin=1in]{geometry}
\usepackage{amsmath,amssymb}
\usepackage{natbib}
\usepackage{graphicx}
\usepackage{booktabs}
\usepackage{xcolor}
\usepackage{hyperref}
\usepackage{cleveref}
\usepackage{tikz}
\usetikzlibrary{arrows.meta, positioning}

\newcommand{\GMSL}{\text{GMSL}}
\newcommand{\GMST}{\text{GMST}}
\newcommand{\CDW}{\text{CDW}}
\newcommand{\ENSO}{\text{ENSO}}

\title{A Hierarchical Bayesian Framework for\\
  Self-Consistent Sea Level Rise Forecasting}

\author{B.\ Minchew}
\date{Working Draft --- \today}

\begin{document}
\maketitle

% ====================================================================
% TARGET: JAMES, Journal of Climate, or Earth's Future
% LENGTH: ~8000-10000 words (full-length)
% FIGURES: 8-12
% TIMELINE: Submission 3-6 months after Papers A and B
% ====================================================================

\begin{abstract}
% ~250 words
We present a hierarchical Bayesian framework for probabilistic sea level rise projections that integrates observational constraints, process-model structural information, and deep-uncertainty treatment into a single self-consistent system.
The framework operates across four levels: (0) global temperature forcing, (1) total GMSL constrained by the observational rate--temperature relationship, (2) component-level decomposition into seven physically distinct contributions (thermal expansion, glaciers, Greenland, West Antarctic, East Antarctic, Antarctic Peninsula, terrestrial water storage) with component-specific forcing variables and a rate-level budget constraint, and (3) deep-uncertainty modules for marine ice-sheet instability.
Three features distinguish the framework from existing approaches.
First, a rate-level soft budget constraint enforces the sea level budget at every calibration time step, creating cross-component information flow via a collider mechanism that transfers information from well-observed components to poorly observed ones.
Second, the directed acyclic graph (DAG) specification supports causal intervention analysis via do-calculus, enabling counterfactual projections under glacier stabilization or climate engineering scenarios.
Third, learned Kennedy \& O'Hagan model-inadequacy terms at each component level prevent the overconfidence that arises from treating parametric models as truth.
The framework is designed for incremental implementation: it is self-consistent at every stage of completeness, with unresolved components absorbed into a budget-constrained residual.
We demonstrate the framework at Stage 1 (thermosteric decomposition) and Stage 2 (glaciers and Greenland), producing projections that are observationally calibrated, budget-consistent, and have explicit model-inadequacy accounting.
\end{abstract}


% ====================================================================
\section{Introduction}
\label{sec:intro}
% ~800 words
% ====================================================================

% 1. Three categories of existing SLR projection approaches, each with
%    fundamental limitations:
%    - Process models: lack parameter exploration, diverge from obs
%    - Semi-empirical: cannot resolve components, wrong functional form at high T
%    - Expert judgment: lacks formal self-consistency
%
% 2. What is needed: a framework that fuses observational constraints,
%    process-model structure, and expert judgment into a single coherent system.
%
% 3. The hierarchical Bayesian approach provides this. It is not new in
%    geophysics (Berliner 2003, Kennedy & O'Hagan 2001), but has not been
%    applied to the SLR projection problem with the full component decomposition
%    and budget constraint.
%
% 4. Paper A established observationally calibrated projections (the top-level
%    constraint). Paper B demonstrated the budget constraint and WAIS treatment.
%    This paper presents the full framework that unifies these elements and
%    adds the component-level architecture, model-inadequacy terms, DAG
%    specification, and causal intervention capability.
%
% 5. The framework is designed for the Arete Glacier Initiative's need to
%    evaluate real glacier stabilization interventions. The counterfactual
%    capability is not hypothetical.


% ====================================================================
\section{Framework Architecture}
\label{sec:architecture}
% ~2500 words
% ====================================================================

\subsection{Hierarchical structure}
\label{sec:hierarchy}
%
% Four levels:
%
% Level 0: Temperature forcing
%   - SSP → radiative forcing → GMST via FaIR
%   - This is exogenous to the SLR framework
%   - But the framework can accept arbitrary forcing trajectories
%     (including intervention scenarios: SAI, CDR)
%
% Level 1: Total GMSL constraint
%   - DOLS rate-temperature relationship (from Paper A)
%   - This provides the top-level observational anchor
%   - The rate-and-state extension for lagged response
%
% Level 2: Component decomposition
%   - 7 components with component-specific forcing variables
%   - Rate-level budget constraint (Eq: sum of rates = observed total rate)
%   - Component-specific models:
%     * TE: linear in GMST (PC prior on quadratic toward zero)
%     * GL: quadratic in GMST with volume-depletion saturation
%     * GrIS: two-component (SMB linear in GMST + discharge in T_NA)
%     * WAIS: forcing-dependent rate ODE (CDW, with MISI threshold)
%     * EAIS: SMB-dominated (precipitation + discharge from marine sectors)
%     * APIS: atmospheric warming, comparatively well-constrained
%     * LWS: secular (dam/GW) + climate-driven (ENSO)
%
% Level 3: Deep uncertainty (A4 module)
%   - Scenario-mixture model for WAIS MISI (from Paper B)
%   - Extends projection envelope beyond historical constraint
%   - Mixture weights updated via budget residual consistency
%
% Figure 1: Schematic of the 4-level hierarchy.

\subsection{Directed acyclic graph}
\label{sec:dag}
%
% Full DAG specification:
% - Node set (Table 2): forcing, mediating, component rate, constraint, 
%   deep uncertainty, observation nodes
% - Edge set with physical justification for each edge
% - Conditional independence relations (CI-1 through CI-5)
% - Critical distinction: generative vs inferential CI
%   The budget node is a collider. During inference, conditioning on it
%   opens paths between components.
%
% Figure 2: TikZ DAG with all nodes and edges.
%
% Intervention targets for do-calculus:
% - Thwaites stabilization: do(rate_WAIS_Thwaites = f(t))
% - CDW reduction: do(T_CDW = T_CDW - Delta_T)
% - SAI: do(F(t) = F_SAI(t))
% - CDR: do(F(t) = F_CDR(t))

\subsection{Rate-level budget constraint}
\label{sec:budget}
%
% Equation: sum_i rate_i(t) ~ N(rate_total_obs(t), sigma_budget^2(t))
%
% sigma_budget^2 = sigma_obs_total^2 + sum_i sigma_extra_i^2
% where sigma_extra_i are learned model-inadequacy parameters (see below).
% sigma_budget participates in inference (changes at each MCMC step).
%
% Why rate-level, not coefficient-level:
% - AIS responds to CDW, not GMST → no shared polynomial
% - TWS secular component has no temperature dependence
% - Rate-level allows component-specific functional forms
%
% Information flow mechanics:
% - Bottom-up: component data constrain component parameters
% - Top-down: total GMSL constrains the sum (collider mechanism)
% - Cross-component: well-constrained components tighten residual for others
%
% Quantify via posterior correlation Cor[r_i, r_j | H_obs, D]
% vs prior correlation (approximately zero).

\subsection{Model inadequacy}
\label{sec:inadequacy}
%
% Kennedy & O'Hagan (2001) framework applied at each component:
%
% rate_i^obs(t) = rate_i^model(t; theta_i) + delta_i(t) + epsilon_i(t)
%
% Phase 1: scalar sigma_extra_i (white-noise discrepancy)
%   - HalfCauchy prior on sigma_extra_i
%   - Adds one parameter per component
%   - Captures magnitude of model inadequacy
%
% Phase 2 (deferred): GP discrepancy delta_i(t) ~ GP(0, k_i)
%   - Captures temporal structure of inadequacy
%   - Activated when Phase 1 residuals show Durbin-Watson structure
%
% Specific model extensions (not discrepancy terms — forward model improvements):
% - Thermosteric: two-timescale deep-ocean lag (deferred unless residuals demand it)
% - Glaciers: volume-depletion saturation correction (required in Phase 1)

\subsection{Staged implementation}
\label{sec:staged}
%
% The framework is self-consistent at every stage:
%
% Stage 0: Total GMSL-GMST constraint only (DOLS + rate-and-state)
%          All components in residual. Complete (Paper A).
% Stage 1: Thermosteric decomposition. Best-constrained component.
%          Residual = cryospheric + TWS. 
% Stage 2: Glaciers + Greenland. Moderately well-observed.
%          Residual = WAIS + EAIS + APIS + TWS.
% Stage 3: TWS (secular + climate-driven). 
%          Residual ≈ WAIS + EAIS + APIS.
% Stage 4: AIS with deep-uncertainty module.
%          Framework complete. No residual.
%
% Table 3: Implementation roadmap (what, data, status).
%
% At each stage: budget holds, posteriors are valid, projections can be made.
% Adding a component refines but does not invalidate.


% ====================================================================
\section{Demonstration: Stages 1-2}
\label{sec:demo}
% ~2000 words
% ====================================================================

\subsection{Stage 1: Thermosteric decomposition}
%
% Model: rate_thermo(T) = b_thermo * T + c_thermo
%   PC prior on a_thermo toward zero (expect linear response).
%   sigma_extra_thermo ~ HalfCauchy(0, s_thermo).
%
% Data: NOAA/NCEI 0-2000m thermosteric, Argo era (2005+).
%   Pre-Argo: Frederikse thermodynamic (1900-2018).
%
% Budget constraint: rate_residual = rate_total - rate_thermo
%   This residual = glaciers + GrIS + WAIS + EAIS + APIS + LWS.
%
% Results:
% - Posterior on b_thermo (linear sensitivity): value, comparison to DOLS alpha_0
% - sigma_extra_thermo: magnitude of thermosteric model inadequacy
% - Budget residual: time series, comparison to IMBIE + GlaMBIE
%
% Figure 3: Stage 1 results. Top: thermosteric rate (model + data + posterior).
% Bottom: budget residual vs independently observed cryospheric + TWS rates.

\subsection{Stage 2: Glaciers and Greenland}
%
% Glacier model: rate_gl(T) = a_gl * T^2 + b_gl * T + c_gl
%   with volume-depletion saturation: rate_gl * [V_gl(t)/V_gl_0]^p
%   V_gl_0 ~ 0.32 m SLE. Exponent p calibrated with prior from Marzeion/Edwards.
%   sigma_extra_gl ~ HalfCauchy.
%
% GrIS model: rate_GrIS = (b_SMB * T + c_SMB) + (beta_dis * T_NA + c_dis)
%   Two-component: SMB (GMST-forced) + discharge (ocean-forced).
%   sigma_extra_GrIS ~ HalfCauchy.
%   T_NA treated as latent variable pre-Argo (parameterized from GMST with uncertainty).
%
% Data: GlaMBIE (glaciers, 2000+), IMBIE GrIS (1992+), RACMO/MAR SMB.
%
% Budget constraint: rate_residual_2 = rate_total - rate_thermo - rate_gl - rate_GrIS
%   This residual ≈ WAIS + EAIS + APIS + LWS.
%
% Results:
% - Glacier saturation exponent p: posterior value
% - GrIS SMB vs discharge partition: posterior on b_SMB, beta_dis
% - Residual tightening from Stage 1 to Stage 2
%
% Figure 4: Stage 2 results. Glacier + GrIS rates. Budget residual narrows.

\subsection{Projection comparison}
%
% Compare projections at Stage 0, 1, 2 under SSP2-4.5:
% - Stage 0: widest uncertainty (all components in residual)
% - Stage 1: narrower (thermosteric resolved)
% - Stage 2: narrower still (glaciers + GrIS resolved)
%
% The residual at Stage 2 provides a budget-constrained estimate of
% WAIS + EAIS + APIS + LWS. Compare to Paper B's direct analysis.
%
% Figure 5: Projection comparison across stages.
% Show how each stage tightens the projection.

\subsection{Model inadequacy diagnosis}
%
% Report sigma_extra_i for each component at Stages 1 and 2.
% If sigma_extra is large relative to sigma_obs, the parametric model
% is missing structure. Flag for Phase 2 GP discrepancy.
%
% Compute Durbin-Watson statistic on residuals for each component.
% If DW indicates autocorrelation, the scalar sigma_extra is insufficient
% and GP discrepancy should be activated.


% ====================================================================
\section{Causal Intervention Analysis}
\label{sec:interventions}
% ~1000 words
% ====================================================================

% Even at Stage 2 (without full AIS model), the DAG supports interventions
% on the forcing nodes:
%
% Example 1: SAI intervention
%   do(F(t) = F_SAI(t)) → modified T(t) → modified component rates
%   With ENSO exogenous (Phase 1), this does NOT propagate to WAIS via ENSO.
%   Document this limitation.
%
% Example 2: Thwaites stabilization
%   do(rate_WAIS_Thwaites = f_stabilized(t))
%   The budget residual changes → other components' posteriors shift.
%   This shows the integrated effect of stabilizing one glacier system.
%
% Present as: "the framework supports these queries in principle.
% At Stage 2, the AIS components are in the residual, so intervention
% analysis on WAIS requires the Paper B mixture model or Stage 4.
% We demonstrate the intervention engine on the forcing-level queries
% (SAI, CDR) where the full component chain is resolved."
%
% Figure 6: Counterfactual projection under SAI scenario.
% Show: total GMSL under SSP2-4.5, then with SAI reducing forcing after 2040.
% The difference = the SLR avoided by the intervention.


% ====================================================================
\section{Variance Decomposition}
\label{sec:variance}
% ~600 words
% ====================================================================

% Nested law of total variance decomposition:
% Total Var = E_S[Var_theta,k[H|S]] + Var_S[E_theta,k[H|S]]
%
% Within-scenario:
% Var_theta,k[H|S] = E_k[Var_theta[H|S,k]] + Var_k[E_theta[H|S,k]]
%
% Cross-component covariance C_cross(t).
%
% R_sep diagnostic for separability of parameter and scenario uncertainty.
%
% Monte Carlo computation (exact up to MC error).
%
% Figure 7: Variance decomposition stacked area plot.
% 7 components + interaction + scenario spread.
% Show WAIS dominating at long horizons (connects to Paper B's result).


% ====================================================================
\section{Relationship to Existing Work}
\label{sec:existing}
% ~500 words
% ====================================================================

% Three primary differentiators from FACTS:
% 1. Rate-level budget constraint (cross-component regularization)
% 2. Causal intervention capability (do-calculus / graph surgery)
% 3. Joint historical calibration (all components calibrated simultaneously
%    against the same observations, not independently)
%
% Positioned as complementary, not adversarial.
%
% Comparison to Kopp et al. (2014, 2017): similar component decomposition
% but without the global budget constraint from rate-temperature calibration.
%
% Comparison to semi-empirical models: the framework generalizes Vermeer &
% Rahmstorf (2009) by decomposing the aggregate relationship into components.


% ====================================================================
\section{Discussion}
\label{sec:discussion}
% ~600 words
% ====================================================================

% Open issues and deferred items:
% - Phase 2 extensions: GP discrepancy, ENSO endogeneity, full AIS state model
% - GRACE cross-product error correlations
% - Paleo data integration
% - Value-of-information analysis for observational investments
%
% The framework is designed to accommodate these without architectural changes.
% Each extension adds detail within the existing hierarchy, not new hierarchy.


% ====================================================================
\section{Conclusions}
\label{sec:conclusions}
% ~300 words
% ====================================================================

% 1. The framework unifies observational calibration, component decomposition,
%    budget constraint, model inadequacy, and deep uncertainty into one system.
% 2. It is self-consistent at every stage of implementation.
% 3. The budget constraint provides cross-component information flow that
%    existing frameworks do not exploit.
% 4. Causal intervention analysis is a structural capability, not an add-on.
% 5. Papers A and B are special cases: Paper A is Stage 0, Paper B is the
%    budget constraint + WAIS deep uncertainty. This paper presents the
%    architecture that unifies them.


% ====================================================================
% SUPPLEMENTARY MATERIALS
% ====================================================================
%
% SM1: Full DAG specification (node set, edge set, CI relations)
% SM2: Rate-level budget constraint derivation
% SM3: Kennedy & O'Hagan model inadequacy: full formulation
% SM4: MCMC specification (sampler choice, tuning, convergence diagnostics)
% SM5: Prior specifications for all parameters (table)
% SM6: Component data details (temporal coverage, processing)
% SM7: Phase 2 design notes (GP discrepancy, ENSO endogeneity)
%
% Fig S1: Full TikZ DAG (larger version with all annotations)
% Fig S2: Prior sensitivity analysis
% Fig S3: MCMC convergence diagnostics (trace plots, R-hat, ESS)
% Fig S4: Posterior predictive checks (model vs observed rates, each component)
% Fig S5: Stage 0 → 1 → 2 residual evolution
% Fig S6: sigma_extra posteriors for each component
% Fig S7: Durbin-Watson diagnostics on component residuals
%
% Table S1: Full prior specification table
% Table S2: Posterior summary statistics (all parameters)
% Table S3: MCMC convergence diagnostics
% Table S4: Projection tables (all stages, all SSPs, decadal)

% MAIN FIGURES (8):
% Fig 1: Hierarchical structure schematic (4 levels)
% Fig 2: DAG (TikZ)
% Fig 3: Stage 1 results (thermosteric + residual)
% Fig 4: Stage 2 results (glaciers + GrIS + narrowed residual)
% Fig 5: Projection comparison across stages
% Fig 6: Counterfactual projection under SAI
% Fig 7: Variance decomposition
% Fig 8: Posterior correlation matrix (7x7)

% TABLES (4):
% Table 1: Component decomposition (7 components, forcing, data, model form)
% Table 2: DAG node set
% Table 3: Implementation roadmap
% Table 4: Projection summary (Stages 0-2, all SSPs)


\bibliographystyle{agsm}
\bibliography{slr_references}

\end{document}
